% !TEX root = ../algebra_02.tex

\section{Расширения полей}

\begin{Definition}{Расширение поля}{}
	Пусть $K$ --- поле. $L$ называется \textbf{расширением} $K$, если $L$ --- поле и $K$ является подполем $L$. Обозначается как $L/K$.
\end{Definition}

\begin{Example}{}{}
	$\mathbb{R}$ --- расширение $\mathbb{Q}$. $\mathbb{C}$ --- расширение $\mathbb{R}$.
\end{Example}

\begin{Remark}{}{}
	Формально поле комплексных чисел можно задать как $\mathbb{C}' = \mathbb{R}[x]/(x^2+1)$, где $\bar{x} = a$. 
	Отождествляя подмножество $\mathbb{R}' := \{ \bar{a} \mid c \in \mathbb{R} \}$ с полем $\mathbb{R}$, получают стандартное поле комплексных чисел: $\mathbb{C} := (\mathbb{C}' \setminus \mathbb{R}') \cup \mathbb{R}$.
\end{Remark}

\subsection{Векторное пространство над подполем}

Любое расширение $L$ поля $K$ можно рассматривать как линейное (векторное) пространство над $K$. 
Операция умножения задана изначально как $L \times L \to L$, $(x,y) \mapsto x \cdot y$. Ограничив её до $K \times L \to L$, получаем операцию умножения вектора из $L$ на скаляр из $K$.

\begin{Definition}{Конечное расширение}{}
	Расширение $L/K$ называется \textbf{конечным}, если $L$ конечномерно как векторное пространство над $K$.
	Размерность этого пространства $\dim_K L$ называется \textbf{степенью расширения} и обозначается $[L:K]$.
\end{Definition}

\begin{Example}{}{}
	$[\mathbb{C}:\mathbb{R}] = 2$. Базисом является множество $\{1, i\}$.
\end{Example}

\begin{Statement}{}{}
	Расширение $L/K$ конечно тогда и только тогда, когда $[L:K] < \infty$.
\end{Statement}

\begin{Theorem}{О мультипликативности степени}{}
	Пусть даны конечные расширения $L/K$ и $K/F$. Тогда расширение $L/F$ также конечно, и выполняется равенство:
	\[ [L:F] = [L:K] \cdot [K:F] \]
\end{Theorem}

\begin{proof}
	Пусть $e_1, \dots, e_m$ --- базис $K$ над $F$, а $\delta_1, \dots, \delta_n$ --- базис $L$ над $K$.
	Докажем, что система элементов $\{ e_i \delta_j \mid i=1\dots m, j=1\dots n \}$ является базисом $L/F$.
	
	1. \textbf{Порождение:} Пусть $a \in L$. Так как $\delta_j$ образуют базис $L/K$, элемент $a$ можно разложить:
	\[ a = \sum_{j=1}^n b_j \delta_j, \quad b_j \in K \]
	Так как $e_i$ образуют базис $K/F$, каждый коэффициент $b_j$ можно разложить:
	\[ b_j = \sum_{i=1}^m c_{ij} e_i, \quad c_{ij} \in F \]
	Подставляя, получаем:
	\[ a = \sum_{j=1}^n \left( \sum_{i=1}^m c_{ij} e_i \right) \delta_j = \sum_{i,j} c_{ij} e_i \delta_j \]
	Таким образом, система $\{ e_i \delta_j \}$ порождает $L$ над $F$.
	
	2. \textbf{Линейная независимость:} Пусть некоторая линейная комбинация равна нулю:
	\[ \sum_{i,j} c_{ij} e_i \delta_j = 0, \quad c_{ij} \in F \]
	Сгруппируем слагаемые при $\delta_j$:
	\[ \sum_{j=1}^n \left( \sum_{i=1}^m c_{ij} e_i \right) \delta_j = 0 \]
	Выражение в скобках является элементом поля $K$. Так как элементы $\delta_j$ линейно независимы над $K$, все коэффициенты при них равны нулю:
	\[ \sum_{i=1}^m c_{ij} e_i = 0 \quad \text{для всех } j=1\dots n \]
	Так как элементы $e_i$ линейно независимы над $F$, отсюда следует, что:
	\[ c_{ij} = 0 \quad \text{для всех } i, j \]
	Следовательно, система линейно независима и образует базис.
\end{proof}

\begin{Remark}{}{}
	Если $F \subset K \subset L$ и расширение $L/F$ конечно, то расширения $L/K$ и $K/F$ также конечны.
\end{Remark}

\subsection{Алгебраические и трансцендентные элементы}

Пусть дано расширение $L/K$ и элемент $a \in L$.

\begin{Definition}{Алгебраический и трансцендентный элементы}{}
	Элемент $a$ называется \textbf{алгебраическим} над $K$, если существует ненулевой многочлен $f \in K[x]$ такой, что $f(a) = 0$.
	В противном случае элемент $a$ называется \textbf{трансцендентным} над $K$.
\end{Definition}

Если элемент $a$ алгебраичен над $K$, рассмотрим множество всех многочленов из $K[x]$, зануляющих $a$:
\[ I = \{ f \in K[x] \mid f(a) = 0 \} \]
Множество $I$ является идеалом в кольце $K[x]$. Так как $K[x]$ --- кольцо главных идеалов, идеал $I$ порождается одним элементом: $I = (f_0)$.

\begin{Definition}{Минимальный многочлен}{}
	Многочлен $f_0$, порождающий идеал $I$, называется \textbf{минимальным многочленом} элемента $a$ над $K$.
	Он определён однозначно с точностью до умножения на обратимый элемент $\varepsilon \in K^*$. Обозначается $\operatorname{Irr}_K(a)$.
\end{Definition}

\begin{Lemma}{Неприводимость минимального многочлена}{}
	Минимальный многочлен $\operatorname{Irr}_K(a)$ неприводим.
\end{Lemma}

\begin{proof}
	Обозначим $f_0 := \operatorname{Irr}_K(a)$. Предположим, что $f_0 = g \cdot h$.
	Так как $f_0(a) = 0$, то $g(a) \cdot h(a) = 0$.
	Поскольку мы работаем в поле, делителей нуля нет, следовательно, либо $g(a) = 0$, либо $h(a) = 0$.
	Это означает, что либо $f_0 \mid g$, либо $f_0 \mid h$.
	Так как степени $g$ и $h$ не превосходят степени $f_0$, это возможно только в том случае, когда $g \sim f_0$ или $h \sim f_0$. Значит, многочлен $f_0$ неприводим.
\end{proof}

\begin{Definition}{Степень алгебраического элемента}{}
	Если элемент $a$ алгебраичен над $K$, то \textbf{степенью} элемента $a$ называется степень его минимального многочлена: $\deg(\operatorname{Irr}_K(a))$.
\end{Definition}

\begin{Example}{}{}
	\begin{enumerate}
		\item $\deg_{\mathbb{R}} i = 2$, так как $\operatorname{Irr}_{\mathbb{R}}(i) = x^2 + 1$.
		\item $\deg_{\mathbb{Q}} \sqrt[3]{2} = 3$, так как $\operatorname{Irr}_{\mathbb{Q}}(\sqrt[3]{2}) = x^3 - 2$.
	\end{enumerate}
\end{Example}

\begin{Proposition}{}{}
	Пусть $[L:K] = n < \infty$. Тогда любой элемент $a \in L$ является алгебраическим над $K$, и его степень не превосходит $n$.
\end{Proposition}

\begin{proof}
	Рассмотрим набор из $n+1$ элементов: $1, a, a^2, \dots, a^n$.
	Так как размерность пространства $\dim_K L = n$, этот набор линейно зависим.
	Следовательно, существуют коэффициенты $\alpha_0, \alpha_1, \dots, \alpha_n \in K$, не все равные нулю, такие что:
	\[ \alpha_0 + \alpha_1 a + \dots + \alpha_n a^n = 0 \]
	Рассмотрим многочлен $f(x) := \alpha_n x^n + \dots + \alpha_1 x + \alpha_0$. Так как не все $\alpha_i$ равны нулю, $f \ne 0$.
	При этом $f(a) = 0$, что означает, что $a$ --- алгебраический элемент над $K$.
	Так как $f(a) = 0$, многочлен $f$ принадлежит идеалу, порожденному $\operatorname{Irr}_K(a)$.
	Следовательно, $\operatorname{Irr}_K(a)$ делит $f$, откуда получаем:
	\[ \deg(\operatorname{Irr}_K(a)) \le \deg(f) \le n \]
\end{proof}

\begin{Remark}{}{}
	$\deg_K a = 1 \iff a \in K$.
\end{Remark}

\subsection{Алгебраические расширения и присоединение элементов}

\begin{Definition}{Алгебраическое расширение}{}
	Расширение $L/K$ называется \textbf{алгебраическим}, если все элементы поля $L$ алгебраичны над $K$.
	В противном случае расширение $L/K$ называется \textbf{трансцендентным}.
\end{Definition}

Пусть $L/K$ --- расширение, и $P \subset L$ --- некоторое подмножество.
Определим $K(P)$ как пересечение всех промежуточных полей $F$, содержащих $K$ и $P$:
\[ K(P) := \bigcap_{\substack{K \subset F \subset L \\ F \text{ --- поле} \\ P \subset F}} F \]
Очевидно, что $K(P)$ является полем. Оно называется полем, полученным из $K$ \textbf{присоединением элементов} из множества $P$.
Если множество элементов конечно ($P = \{x_1, \dots, x_n\}$), используют обозначение $K(x_1, \dots, x_n)$.

\begin{Definition}{Конечно порождённое и простое расширения}{}
	Расширение $L/K$ называется \textbf{конечно порождённым}, если существуют элементы $x_1, \dots, x_n \in L$ такие, что $L = K(x_1, \dots, x_n)$.
	Расширение $L/K$ называется \textbf{простым}, если оно порождается одним элементом: $L = K(x)$ для некоторого $x \in L$.
\end{Definition}

\begin{Remark}{}{}
	Если элемент $a$ трансцендентен над $K$, то расширение $L = K(a)$ называется простым трансцендентным. Если $a$ алгебраичен, то $L = K(a)$ называется простым алгебраическим расширением.
\end{Remark}

\begin{Theorem}{Критерий конечности расширения}{}
	Для расширения $L/K$ следующие условия эквивалентны:
	\begin{enumerate}
		\item $L/K$ конечно;
		\item $L/K$ алгебраическое и конечно порождённое.
	\end{enumerate}
\end{Theorem}

\subsection{Строение простого алгебраического расширения}

\begin{Theorem}{Об изоморфизме простого расширения}{}
	Пусть $L = K(x)$ --- простое расширение, где элемент $x \in L$ алгебраичен над $K$. Пусть $f_0 = \operatorname{Irr}_K(x)$ --- его минимальный многочлен, а $\deg f_0 = d$. 
	Тогда существует изоморфизм полей:
	\[ K[t]/(f_0) \xrightarrow{\sim} L \]
	При этом элементы смежных классов вида $\alpha_0 + \alpha_1 t + \dots + \alpha_{d-1} t^{d-1}$ переходят в элементы $\alpha_0 + \alpha_1 x + \dots + \alpha_{d-1} x^{d-1} \in L$.
\end{Theorem}

\begin{proof}
	Рассмотрим отображение вычисления многочлена в точке $x$:
	\[ \varphi: K[t] \to L, \quad \sum_{i=0}^n \alpha_i t^i \mapsto \sum_{i=0}^n \alpha_i x^i \]
	Очевидно, что $\varphi$ является гомоморфизмом колец.
	Найдем ядро этого гомоморфизма. По определению, многочлен $f(t)$ зануляется в точке $x$ тогда и только тогда, когда он делится на минимальный многочлен $f_0$. Следовательно, $\ker \varphi = (f_0)$.
	
	По теореме о гомоморфизме колец, индуцированное отображение $\tilde{\varphi}$ задаёт изоморфизм:
	\[ K[t]/(f_0) \xrightarrow{\sim} \operatorname{Im} \varphi \]
	Поскольку минимальный многочлен $f_0$ неприводим, идеал $(f_0)$ максимален, а значит, факторкольцо $K[t]/(f_0)$ является полем. Следовательно, образ $\operatorname{Im} \varphi$ также является полем.
	
	Заметим, что:
	\begin{itemize}
		\item $\operatorname{Im} \varphi \supset K$, так как для скаляров $\tilde{\varphi}(\alpha) = \alpha$;
		\item $\operatorname{Im} \varphi \ni x$, так как $\tilde{\varphi}(\bar{t}) = x$;
		\item $\operatorname{Im} \varphi \subset L$ по построению.
	\end{itemize}
	Так как $L = K(x)$ --- наименьшее поле, содержащее $K$ и $x$, а $\operatorname{Im} \varphi$ --- поле с такими же свойствами, получаем $\operatorname{Im} \varphi = L$. 
\end{proof}

\begin{Definition}{$K$-изоморфизм расширений}{}
	Пусть даны расширения $L/K$ и $L'/K$. Изоморфизм полей $\varphi: L \to L'$ называется \textbf{$K$-изоморфизмом расширений}, если он оставляет все элементы базового поля $K$ на месте:
	\[ \forall a \in K: \varphi(a) = a \]
\end{Definition}

Таким образом, доказанная теорема показывает, что простое алгебраическое расширение $K(x)$ единственно с точностью до $K$-изоморфизма и полностью определяется минимальным многочленом элемента $x$.

\begin{Example}{}{}
	Рассмотрим многочлен $f_0 = x^3 - 2 \in \mathbb{Q}[x]$. Он неприводим над $\mathbb{Q}$.
	Корнями этого многочлена в $\mathbb{C}$ являются числа $\sqrt[3]{2}$, $\sqrt[3]{2}\omega$ и $\sqrt[3]{2}\omega^2$, где $\omega = e^{2\pi i / 3}$.
	Поскольку все эти элементы имеют один и тот же минимальный многочлен $f_0$, порожденные ими простые расширения изоморфны:
	\[ \mathbb{Q}(\sqrt[3]{2}) \cong \mathbb{Q}(\sqrt[3]{2}\omega) \cong \mathbb{Q}(\sqrt[3]{2}\omega^2) \cong \mathbb{Q}[x]/(x^3 - 2) \]
\end{Example}

\begin{Theorem}{О существовании простого расширения}{}
	Пусть $f_0 \in K[t]$ --- неприводимый многочлен. Тогда существует расширение $L/K$ и элемент $x \in L$ такие, что $L = K(x)$ и $\operatorname{Irr}_K(x) = f_0$.
\end{Theorem}

\begin{proof}
	Сконструируем поле $L = K[t]/(f_0)$. Мы можем рассматривать $L$ как расширение $K$, отождествляя каждый элемент $\alpha \in K$ с классом смежности скалярного многочлена $\bar{\alpha} \in K[t]/(f_0)$.
	Обозначим класс смежности переменной $t$ за $x = \bar{t}$. Тогда любой элемент из $L$ выражается как многочлен от $x$, то есть $L = K(x)$.
	
	Покажем, что $x$ --- корень $f_0$. Если $f_0(t) = \sum \alpha_i t^i$, то:
	\[ f_0(x) = \sum \alpha_i \bar{t}^i = \overline{\sum \alpha_i t^i} = \overline{f_0(t)} = 0 \]
	Таким образом, $f_0 \in \operatorname{Ann}(x)$. При этом для любого многочлена $f \in K[t]$ выполняется:
	\[ f(x) = 0 \iff \overline{f(t)} = 0 \iff f(t) \in (f_0) \iff f_0 \mid f \]
	Следовательно, $f_0$ --- многочлен наименьшей степени, зануляющий $x$, то есть $\operatorname{Irr}_K(x) = f_0$.
\end{proof}

\begin{Corollary}{}{}
	Если элемент $x$ алгебраичен над $K$, то степень порожденного им простого расширения равна степени его минимального многочлена:
	\[ [K(x):K] = \deg \operatorname{Irr}_K(x) \]
\end{Corollary}

\begin{Corollary}{}{}
	Пусть $L/K$ --- конечное расширение, а $x \in L$. Тогда степень алгебраического элемента $x$ делит степень расширения:
	\[ \deg_K x \mid [L:K] \]
\end{Corollary}

\begin{proof}
	Рассмотрим башню полей $K \subset K(x) \subset L$.
	По теореме о мультипликативности степени конечных расширений:
	\[ [L:K] = [L:K(x)] \cdot [K(x):K] \]
	Подставляя $[K(x):K] = \deg_K x$, получаем требуемое утверждение.
\end{proof}

\subsection{Конечно порожденные алгебраические расширения}

\begin{Theorem}{Критерий конечности расширения (доказательство)}{}
	Для расширения $L/K$ следующие условия эквивалентны:
	\begin{enumerate}
		\item $L/K$ конечно;
		\item $L/K$ конечно порождено и алгебраично.
	\end{enumerate}
\end{Theorem}

\begin{proof}
	\textbf{$(1 \implies 2)$:} Пусть расширение $L/K$ конечно. Тогда оно имеет конечный базис $a_1, \dots, a_n$ как векторное пространство над $K$. 
	Рассмотрим поле $K_1 = K(a_1, \dots, a_n)$. Так как $K_1$ содержит все $a_i$ и поле $K$, оно обязано содержать все их линейные комбинации:
	\[ \forall \alpha_1, \dots, \alpha_n \in K: \alpha_1 a_1 + \dots + \alpha_n a_n \in K_1 \]
	Но эти линейные комбинации покрывают всё пространство $L$, следовательно, $L = K_1 = K(a_1, \dots, a_n)$. Таким образом, расширение конечно порождено. Алгебраичность любого конечного расширения была доказана ранее.
	
	\textbf{$(2 \implies 1)$:} Пусть $L = K(a_1, \dots, a_n)$, где все $a_i$ алгебраичны над $K$. Докажем конечность расширения индукцией по количеству порождающих элементов $j$.
	Обозначим $L_j := K(a_1, \dots, a_j)$ для $j = 0, 1, \dots, n$.
	\begin{itemize}
		\item \textit{База:} При $j=0$ имеем $L_0 = K$, и степень расширения $[K:K] = 1 < \infty$.
		\item \textit{Шаг:} Пусть для $j-1$ доказано, что $[L_{j-1}:K] < \infty$.
		Рассмотрим поле $L_j = L_{j-1}(a_j)$. Элемент $a_j$ алгебраичен над $K$, а значит, подавно алгебраичен и над более широким полем $L_{j-1}$. 
		Поэтому степень простого расширения $[L_j:L_{j-1}] \le \deg_K(a_j) < \infty$.
		По теореме о башне конечных расширений:
		\[ [L_j:K] = [L_j:L_{j-1}] \cdot [L_{j-1}:K] < \infty \]
	\end{itemize}
	Применяя индукционный переход вплоть до $j=n$, получаем $[L_n:K] = [L:K] < \infty$.
\end{proof}

\begin{Proposition}{Поле алгебраических чисел}{}
	Пусть дано произвольное расширение $L/K$. Множество всех элементов поля $L$, алгебраических над $K$:
	\[ M = \{ x \in L \mid x \text{ алгебраичен над } K \} \]
	является подполем в $L$.
\end{Proposition}

\begin{proof}
	Необходимо проверить, что множество $M$ содержит $0$ и $1$, а также замкнуто относительно сложения, умножения и взятия обратного элемента.
	
	Очевидно, что $0, 1 \in M$, так как элементы базового поля $K$ являются корнями многочленов первой степени (алгебраичны).
	Пусть $x, y \in M$, причем $x, y \neq 0$. Рассмотрим промежуточное поле $K_1 = K(x, y)$. 
	Так как $x$ и $y$ алгебраичны над $K$, расширение $K_1/K$ является конечно порожденным алгебраическим. 
	По критерию конечности, $[K_1:K] < \infty$.
	
	Любое конечное расширение является алгебраическим, поэтому все элементы поля $K_1$ алгебраичны над $K$. 
	Так как элементы $x+y$, $x \cdot y$ и $x^{-1}$ принадлежат $K_1$, они также алгебраичны над $K$ и, следовательно, лежат в $M$. Это доказывает, что $M$ --- поле.
\end{proof}

\subsection{Поле разложения многочлена}

\begin{Definition}{Поле разложения}{}
	Пусть $K$ --- поле, $f \in K[x]$ --- ненулевой многочлен.
	Расширение $L/K$ называется \textbf{полем разложения} многочлена $f$, если выполняются два условия:
	\begin{enumerate}
		\item В кольце многочленов $L[x]$ многочлен $f$ раскладывается на линейные множители:
		\[ f(x) = c(x-a_1)\dots(x-a_n), \quad \text{где } a_i \in L \]
		\item Поле $L$ порождается над $K$ всеми корнями этого многочлена:
		\[ L = K(a_1, \dots, a_n) \]
	\end{enumerate}
\end{Definition}